\begin{tikzpicture}[
  x=1mm,y=1mm,figure picture,
  mtinp/.style={figure node,minimum height=6.5mm,minimum width=13mm,
    inner xsep=1pt,draw=FigureInput,fill=FigureInputFill,
    font=\figurefont\scriptsize},
  mtenc/.style={figure node,minimum height=6.5mm,minimum width=13mm,
    inner xsep=1pt,draw=FigureModel,fill=FigureModelFill,
    font=\figurefont\scriptsize},
  mtwide/.style={mtenc,minimum width=30mm},
  mtrep/.style={figure node,minimum height=6.5mm,minimum width=13mm,
    inner xsep=1pt,font=\figurefont\scriptsize},
  mtout/.style={figure node,minimum height=7mm,minimum width=30mm,
    inner xsep=1pt,draw=FigureOutput,fill=FigureOutputFill,
    font=\figurefont\scriptsize},
  mtnote/.style={font=\figurefont\scriptsize,text=FigureMuted,align=center},
  mtfl/.style={figure flow}
]
  % 四个面板：左下 x、宽 35、高 64
  \figurePanel{0}{0}{35}{64}{表示型}{FigureInputFill}
  \figurePanel{38}{0}{35}{64}{交互型}{FigureInputFill}
  \figurePanel{76}{0}{35}{64}{后期交互}{FigureInputFill}
  \figurePanel{114}{0}{35}{64}{多向量（Poly）}{FigureInputFill}

  % 共同的输入行
  \foreach \x in {0,38,76,114}{
    \node[mtinp] (qa\x) at (\x+9.5,10) {查询$a$};
    \node[mtinp] (cb\x) at (\x+25.5,10) {候选$b$};
  }

  % 表示型：两塔独立编码，单向量比较
  \node[mtenc] (e1) at (9.5,22) {编码器};
  \node[mtenc] (e2) at (25.5,22) {编码器};
  \node[mtrep] (u) at (9.5,34) {$\bm u\in\mathbb R^{d}$};
  \node[mtrep] (v) at (25.5,34) {$\bm v\in\mathbb R^{d}$};
  \node[mtout] (s1) at (17.5,46) {$s=\cos(\bm u,\bm v)$};
  \draw[mtfl] (qa0)--(e1); \draw[mtfl] (cb0)--(e2);
  \draw[mtfl] (e1)--(u); \draw[mtfl] (e2)--(v);
  \draw[mtfl] (u)--(s1.south -| u); \draw[mtfl] (v)--(s1.south -| v);

  % 交互型：联合编码
  \node[mtwide] (j) at (55.5,22) {联合编码$[a;b]$};
  \node[mtrep,minimum width=30mm] (h) at (55.5,34) {句对表示$\bm h$};
  \node[mtout] (s2) at (55.5,46) {$s=\bm w^{\top}\bm h$};
  \draw[mtfl] (qa38)--(qa38 |- j.south); \draw[mtfl] (cb38)--(cb38 |- j.south);
  \draw[mtfl] (j)--(h); \draw[mtfl] (h)--(s2);

  % 后期交互：保留词元向量
  \node[mtenc] (l1) at (85.5,22) {编码器};
  \node[mtenc] (l2) at (101.5,22) {编码器};
  \node[mtrep] (qi) at (85.5,34) {$\{\bm q_i\}_{i\le n_q}$};
  \node[mtrep] (dj) at (101.5,34) {$\{\bm d_j\}_{j\le n_d}$};
  \node[mtout] (s3) at (93.5,46) {$\sum_i\max_j\bm q_i^{\top}\bm d_j$};
  \draw[mtfl] (qa76)--(l1); \draw[mtfl] (cb76)--(l2);
  \draw[mtfl] (l1)--(qi); \draw[mtfl] (l2)--(dj);
  \draw[mtfl] (qi)--(s3.south -| qi); \draw[mtfl] (dj)--(s3.south -| dj);

  % 多向量：m 个代码向量汇聚查询
  \node[mtenc] (p1) at (123.5,22) {编码器};
  \node[mtenc] (p2) at (139.5,22) {编码器};
  \node[mtrep] (ym) at (123.5,34) {$\bm y^{1},\ldots,\bm y^{m}$};
  \node[mtrep] (yc) at (139.5,34) {$\bm y_{\mathrm{cand}}$};
  \node[mtout] (s4) at (131.5,46) {按$\bm y_{\mathrm{cand}}$汇聚后点积};
  \draw[mtfl] (qa114)--(p1); \draw[mtfl] (cb114)--(p2);
  \draw[mtfl] (p1)--(ym); \draw[mtfl] (p2)--(yc);
  \draw[mtfl] (ym)--(s4.south -| ym); \draw[mtfl] (yc)--(s4.south -| yc);

  % 预计算说明
  \node[mtnote] at (17.5,3) {候选向量可离线计算};
  \node[mtnote] at (55.5,3) {每个句对重新编码};
  \node[mtnote] at (93.5,3) {离线存词元向量};
  \node[mtnote] at (131.5,3) {离线存单向量};
\end{tikzpicture}
